\section{Results}
\label{sec:results}

\begin{table}[tb]
\centering\small
\caption{Main results, all figures percentages, scored against the curated
reference. The two blocks are \emph{different models} on \emph{different
splits}. Above: the system we submit, against the published baseline, both on
the established split. Below: its sibling, trained on the text-disjoint split
and therefore the only model we can measure without sentence overlap; its
held-out row was measured once, after every decision was frozen. The two
upper rows are single models, so no seed statistics apply to them.}
\label{tab:main}
\begin{tabular}{lrrrrrr}
\toprule
& \multicolumn{2}{c}{Scripted} & \multicolumn{2}{c}{Spontaneous} & \multicolumn{2}{c}{Mean} \\
\cmidrule(lr){2-3}\cmidrule(lr){4-5}\cmidrule(lr){6-7}
System & WER & CER & WER & CER & WER & CER \\
\midrule
\multicolumn{7}{l}{\emph{Submitted model, on \texttt{sharedtask} development (speaker-disjoint)}} \\
\quad Published baseline, converted & 16.08 & 3.77 & 10.23 & 1.39 & 13.15 & 2.58 \\
\quad This work & 9.23 & 3.01 & 9.06 & 1.30 & \textbf{9.14} & 2.16 \\
\addlinespace
\multicolumn{7}{l}{\emph{Sibling model, trained and measured on \texttt{sharedtask\_both}}} \\
\multicolumn{7}{l}{\emph{(speaker- and text-disjoint), mean of three seeds}} \\
\quad on its development set & 4.85 & 0.62 & 9.20 & 1.31 & 7.02 & 0.96 \\
\quad\quad \emph{std.\ dev.\ across seeds} & 0.21 & 0.04 & 0.18 & 0.02 & 0.19 & 0.03 \\
\quad on its held-out set & 10.79 & 2.33 & 3.38 & 0.39 & 7.08 & 1.36 \\
\quad\quad \emph{std.\ dev.\ across seeds} & 0.14 & 0.02 & 0.07 & 0.01 & 0.05 & 0.01 \\
\bottomrule
\end{tabular}
\end{table}

\subsection{How long to train}
\label{sec:plateau}

Fine-tuning from an already-adapted checkpoint needs far less training than the
original 8{,}000 updates. On the text-disjoint split, every configuration
between 250 and 1{,}000 updates and between $1\times10^{-5}$ and
$5\times10^{-5}$ lands at 7.1--7.3 mean WER: below about 200 updates the
curated orthography is not yet learned, beyond about 1{,}000 the model
memorises, and by 8{,}000 it is worse than at 500. The useful conclusion is a
broad plateau, not an optimum.

Reaching that conclusion required a correction worth reporting. A single run at
500 updates and $2\times10^{-5}$ gave 6.96 mean WER, best of every
configuration on all four numbers, and we recorded it as the winner; a second
seed of the identical configuration gave 7.26. Three seeds measure the noise directly:
7.02 mean WER with standard deviation 0.19 on the text-disjoint development
set, and 9.23 with deviation 0.08 on the larger speaker-disjoint one, of which
Table~\ref{tab:main} reports the submitted seed. We treat any difference below
roughly 0.3 mean WER as indistinguishable from the random draw. The contrast between the two sets is easy to overlook when choosing one:
the smaller does not merely measure less precisely, it manufactures differences.

\subsection{Sentence overlap, measured}
\label{sec:overlap}

Section~\ref{sec:newsentences} makes the established split's text overlap a
measurement problem rather than a curiosity. The design holds the evaluation fixed and varies only what the model saw.
Condition~A trains on the text-disjoint split, so the development sentences
never appear. Condition~B adds the recordings of those same sentences made by
\emph{training} speakers, 3.65 hours more audio. Comparing A with B alone would
confound exposure with data volume, so condition~B$_{\text{eq}}$ removes 3.64
hours of randomly chosen A clips, leaving exactly A's 39.94 hours. Every
condition uses the same configuration, 250 updates, and the same development set.


\begin{table}[tb]
\centering\small
\caption{What hearing a development sentence from other speakers is worth.
Same development set and same configuration (250 updates) throughout; only the
training data differs. Two seeds where two figures are given.}
\label{tab:overlap}
\begin{tabular}{llrr}
\toprule
Training set & Dev sentences heard? & Scripted WER & Scripted CER \\
\midrule
A \quad (39.94 h) & no & 5.20 / 4.97 & 0.65 / 0.62 \\
B$_{\text{eq}}$ (39.94 h) & yes, same hours as A & \textbf{3.79 / 4.23} & \textbf{0.48 / 0.55} \\
B \quad (43.58 h) & yes, with 3.65 h more audio & 4.68 & 0.60 \\
\bottomrule
\end{tabular}
\end{table}

The seed ranges of A and B$_{\text{eq}}$ do not overlap -- the worse
B$_{\text{eq}}$ run still beats the better A run -- so at equal training
volume, having heard a sentence read by other speakers is worth 1.1 WER points,
a 21\% relative reduction, and about 19\% of CER.

It is a lower bound: memorisation grows with training, and the original
baseline runs for 8{,}000 updates.

\subsection{Held-out evaluation}
\label{sec:heldout-results}


The lower block of Table~\ref{tab:main} gives the measurement. The means agree
closely, 7.02 against 7.08, which is the result we care about:
after dozens of decisions taken on the development set, the held-out estimate
did not collapse, and seed spread falls to 0.05.

The agreement hides two large opposing movements, both composition rather than
generalisation. Scripted error doubles because the held-out portion contains
one speaker who alone scores 98.67 WER over 59 clips; removing that same
speaker from the development set moves its scripted WER from 9.23 to 5.71.
Spontaneous error falls by two thirds because 46\% of the held-out spontaneous
portion comes from speakers already heard through silver transcriptions in
training. With 77 scripted speakers, a split places few voices in each
partition, so no such figure should be read without its per-speaker breakdown.

% =====================================================================
